\section{Extension Task: Sequential Bayesian Updating}
\label{sec:extension}

\subsection{Choice of Extension and Justification}

Our group selected Extension Task C: Sequential Bayesian Updating. This choice was
motivated by three considerations. First, sequential updating is the most direct
demonstration of the foundational Bayesian principle that the posterior from one
analysis can serve as the prior for the next, illustrating how beliefs are rationally
revised as evidence accumulates. Second, the subsampling strategy adopted in
Section~\ref{sec:introduction} --- $n = 360$ observations drawn with
\texttt{random\_state = 42} --- was deliberately designed to yield four equal
chronological blocks of 90 observations each, approximating semi-annual windows
over the two-year dataset and making sequential updating both natural and
interpretable. Third, the extension required no modelling machinery beyond what was
already developed in Core Task 1, allowing us to focus on conceptual depth rather
than technical complexity.

\subsection{Methodology}

\subsubsection{Data Partitioning}

The 360 subsampled observations are ordered chronologically by \texttt{dteday} and
partitioned into four non-overlapping blocks of equal size via
\texttt{df.iloc[i*90:(i+1)*90]} for $i = 0, 1, 2, 3$.
Table~\ref{tab:chunks} summarises the temporal coverage and descriptive statistics
of each block.

\begin{table}[htbp]
\centering
\caption{Chronological data blocks used for Sequential Bayesian Updating.
$\bar{y}$ denotes mean daily rentals; $r_{yT}$ and $r_{yH}$ denote the
Pearson correlations of \texttt{cnt} with \texttt{temp} and \texttt{hum}
respectively; $r_{TH}$ denotes the correlation between \texttt{temp} and
\texttt{hum}.}
\label{tab:chunks}
\begin{tabular}{cllcccccc}
\hline
Block & Period & $n$ & $\bar{y}$ & $r_{yT}$ & $r_{yH}$ & $r_{TH}$ \\
\hline
1 & 2011-01-01 -- 2011-06-27 & 90 & 2773 & $+$0.889 & $+$0.165 & $+$0.226 \\
2 & 2011-06-28 -- 2011-12-13 & 90 & 3968 & $+$0.417 & $-$0.518 & $-$0.237 \\
3 & 2011-12-16 -- 2012-06-15 & 90 & 4641 & $+$0.821 & $-$0.054 & $+$0.246 \\
4 & 2012-06-16 -- 2012-12-31 & 90 & 6188 & $+$0.563 & $-$0.188 & $+$0.048 \\
\hline
\end{tabular}
\end{table}

Several features of Table~\ref{tab:chunks} are worth noting before interpreting the
sequential posteriors. The raw correlation $r_{yH}$ is positive in Block 1
($+$0.165) and near zero in Block 3 ($-$0.054), in contrast to the overall
negative relationship reported in Section~\ref{sec:introduction}. Furthermore,
$r_{TH}$ is at most 0.246 across all blocks, ruling out severe multicollinearity
between the two predictors as a source of instability. The heterogeneity of these
local correlations across blocks reflects genuine seasonal variation in the
temperature--humidity--rental relationship, a point we return to in
Section~\ref{sec:ext_discussion}.

\subsubsection{Cascading Prior Update via Moment Matching}

Let $p_k(\boldsymbol{\theta})$ denote the posterior obtained after processing
block $k$, and let $\mu_k^{(j)}$ and $s_k^{(j)}$ denote its mean and standard
deviation for parameter $\beta_j$. The prior for block $k+1$ is constructed by
\emph{moment matching}: the posterior summary statistics replace the
hyperparameters of the next block's prior,

\begin{equation}
    \beta_j \mid \text{block } k+1 \;\sim\;
    \mathcal{N}\!\left(\mu_k^{(j)},\; \bigl(s_k^{(j)}\bigr)^2\right),
    \qquad j = 0, 1, 2,
    \label{eq:moment_match_beta}
\end{equation}

\noindent and for the residual standard deviation,

\begin{equation}
    \sigma \mid \text{block } k+1 \;\sim\;
    \operatorname{Half-Normal}\!\left(\mu_k^{(\sigma)}\right).
    \label{eq:moment_match_sigma}
\end{equation}

Block 1 uses the weakly informative priors specified in Section~\ref{sec:prior}
as its starting point:
$\beta_j \sim \mathcal{N}(0,\,10{,}000^2)$ and
$\sigma \sim \operatorname{Half-Normal}(3{,}000)$.
Each block is fitted with the same NUTS configuration as Core Task 1: 4 chains,
1{,}000 warm-up iterations (discarded as burn-in), and 2{,}000 retained draws,
yielding 8{,}000 posterior samples per block.

\subsection{Convergence Diagnostics}

Before interpreting the sequential posteriors, we verify that MCMC sampling
converged reliably for every block. The Gelman--Rubin statistic $\hat{R}$ compares
between-chain variance to within-chain variance across the four parallel chains.
Formally,

\begin{equation}
    \hat{R} = \sqrt{\frac{\hat{V}}{W}},
\end{equation}

\noindent where $W$ is the mean within-chain variance and $\hat{V}$ is a
weighted combination of $W$ and the between-chain variance $B$. When the chains
have all converged to the same target distribution, $\hat{R} \to 1$; values
above 1.01 are conventionally taken as evidence of non-convergence. As reported
in Table~\ref{tab:seq_posterior}, $\hat{R} = 1.000$ for every parameter in
every block, confirming that the four chains reached the same posterior
distribution in all cases.

The effective sample size (ESS) further quantifies sampling quality by
accounting for autocorrelation between successive draws within each chain.
It answers the question: how many independent samples would carry the same
statistical information as the 8{,}000 autocorrelated draws? Across all blocks
and parameters, bulk ESS ranges from 3{,}792 to 6{,}604, well above the
conventional threshold of 400. Together, $\hat{R} = 1.000$ and high ESS
confirm that the sequential posteriors are reliable estimates and not artefacts
of poor sampling.

\subsection{Results}

Table~\ref{tab:seq_posterior} reports the full posterior summaries for all four
parameters across the four blocks.

\begin{table}[htbp]
\centering
\caption{Sequential posterior summaries. Mean, standard deviation (SD), and
95\% highest-density credible interval (CrI) for each parameter across the four
blocks. All $\hat{R} = 1.000$; ESS values are given in parentheses.}
\label{tab:seq_posterior}
\small
\begin{tabular}{clccccc}
\hline
Block & Parameter & Mean & SD & 95\% CrI lower & 95\% CrI upper & ESS \\
\hline
\multirow{4}{*}{1} 
  & $\beta_0$ &   110.7 & 274.6 & $-$436.5 &   636.1 & 3792 \\
  & $\beta_1$ &  7048.0 & 403.7 &  6227.1  &  7800.2 & 4555 \\
  & $\beta_2$ &  $-$313.0 & 427.1 & $-$1141.5 &   539.1 & 3901 \\
  & $\sigma$  &   652.3 &  51.4 &   554.3  &   753.7 & 5179 \\
\hline
\multirow{4}{*}{2}
  & $\beta_0$ &   462.3 & 220.5 &    40.9  &   902.8 & 4919 \\
  & $\beta_1$ &  6343.3 & 322.9 &  5683.0  &  6962.2 & 5145 \\
  & $\beta_2$ &  $-$248.3 & 299.4 &  $-$836.0 &   338.7 & 5430 \\
  & $\sigma$  &  1127.3 &  86.4 &   965.5  &  1301.5 & 4959 \\
\hline
\multirow{4}{*}{3}
  & $\beta_0$ &  1214.6 & 165.3 &   910.1  &  1547.6 & 4419 \\
  & $\beta_1$ &  7405.7 & 288.1 &  6838.4  &  7963.9 & 5060 \\
  & $\beta_2$ &  $+$180.4 & 250.0 &  $-$297.7 &   686.2 & 4903 \\
  & $\sigma$  &  1088.4 &  88.4 &   924.5  &  1267.4 & 5544 \\
\hline
\multirow{4}{*}{4}
  & $\beta_0$ &  1490.6 & 141.3 &  1209.0  &  1761.3 & 6420 \\
  & $\beta_1$ &  7606.9 & 241.6 &  7135.4  &  8076.6 & 6030 \\
  & $\beta_2$ &  $+$438.5 & 216.3 &    14.1  &   853.0 & 6286 \\
  & $\sigma$  &  1497.1 & 113.4 &  1279.6  &  1720.1 & 6604 \\
\hline
\end{tabular}
\end{table}

\subsubsection{Cascade Plot}

Figure~\ref{fig:cascade} displays the posterior means and 95\% credible intervals
for all four parameters across the four sequential blocks.

\begin{figure}[htbp]
    \centering
    \includegraphics[width=\textwidth]{figures/cascade_plot.png}
    \caption{Cascade plot showing the evolution of posterior means (points) and
    95\% credible intervals (vertical lines) across the four sequential blocks.
    Darker shading corresponds to later blocks. The progressive narrowing of the
    intervals for $\beta_1$ and $\sigma$ reflects the accumulation of information,
    while $\beta_2$ remains uncertain across all blocks due to small-sample and
    seasonal effects discussed in Section~\ref{sec:ext_discussion}.}
    \label{fig:cascade}
\end{figure}

\subsubsection{Density Evolution Plot}

Figure~\ref{fig:density} presents the posterior density curves for all four
parameters overlaid across the four blocks. This representation makes the
distributional narrowing more explicit than the interval plot alone.

\begin{figure}[htbp]
    \centering
    \includegraphics[width=\textwidth]{figures/density_evolution.png}
    \caption{Posterior density evolution across the four sequential blocks.
    Lighter shading corresponds to earlier blocks; darker shading to later ones.
    The densities for $\beta_0$ and $\beta_1$ progressively concentrate around
    stable values, demonstrating Bayesian learning. The $\beta_2$ densities shift
    substantially in location across blocks, reflecting the instability documented
    in Section~\ref{sec:ext_discussion}.}
    \label{fig:density}
\end{figure}

\subsection{Discussion}
\label{sec:ext_discussion}

\subsubsection{Temperature Coefficient $\beta_1$: Convergence with Increasing Data}

The temperature coefficient $\beta_1$ exhibits the clearest Bayesian learning
behaviour. Its posterior standard deviation decreases monotonically from 403.7
in Block 1 to 241.6 in Block 4, a reduction of 40\%, and its 95\% credible
interval narrows from a width of 1{,}573 to 941. The posterior mean fluctuates
across blocks (7{,}048 $\to$ 6{,}343 $\to$ 7{,}406 $\to$ 7{,}607) but remains
entirely positive in all four blocks, confirming that the positive effect of
temperature on rentals is a consistent finding at the local level.

The Block 4 estimate of 7{,}607 is somewhat higher than the full-data posterior
mean of approximately 6{,}784 reported in Section~\ref{sec:posterior}. This
discrepancy is explained by the composition of Block 4 (June--December 2012),
which has the highest mean daily rentals of any block ($\bar{y} = 6{,}188$),
more than double that of Block 1 ($\bar{y} = 2{,}773$). Within this high-demand
summer--autumn period, temperature variation explains a larger share of rental
variation, inflating the local estimate of $\beta_1$ relative to the full-year
average.

\subsubsection{Humidity Coefficient $\beta_2$: Small-Sample Noise and Seasonal Effects}

The humidity coefficient $\beta_2$ presents a more complex picture. Its posterior
mean changes sign from negative in Blocks 1--2 ($-$313, $-$248) to positive in
Blocks 3--4 ($+$180, $+$438), and the full-data estimate of approximately
$-$2{,}350 from Section~\ref{sec:posterior} is not recovered at any block.
Two complementary explanations account for this behaviour.

\textbf{Small-sample noise.} The 95\% credible intervals for $\beta_2$ span zero
in Blocks 1, 2, and 3 ($[-1142, 539]$, $[-836, 339]$, $[-298, 686]$ respectively),
indicating that the humidity effect is statistically unresolved within any single
90-observation window. Only Block 4 yields a credible interval entirely above
zero ($[14, 853]$), and even this result contradicts the direction established
by the full dataset. The contrast with $\beta_1$, whose credible interval is
entirely positive in every block, illustrates that 90 observations are sufficient
to identify the dominant temperature effect but not the weaker and more
variable humidity effect.

\textbf{Seasonal variation.} Table~\ref{tab:chunks} reveals genuine heterogeneity
in the raw correlation between humidity and rentals across blocks
($r_{yH}$: $+$0.165, $-$0.518, $-$0.054, $-$0.188). Crucially, multicollinearity
between \texttt{temp} and \texttt{hum} can be ruled out as a confounding factor:
the between-predictor correlation $r_{TH}$ reaches at most 0.246, which is
insufficient to cause serious instability in the regression coefficients. Instead,
the sign reversal of $r_{yH}$ in Block 1 reflects a genuine feature of the
winter--spring period: rising temperatures and rising humidity co-occur as spring
approaches, so humidity is locally positively associated with the higher rental
volumes of warmer months. This local structure differs from the global
temperature-adjusted humidity effect captured by $\beta_2$ in the full dataset.

The propagation of moment-matched priors compounds this instability: the Block 1
estimate of $\beta_2 = -313$ is passed as the prior mean to Block 2,
anchoring its posterior near $-$248 rather than the true value of $-$2{,}350.
Each subsequent block inherits a prior already displaced from the global estimate,
and the small within-block samples lack the power to correct the accumulated
displacement.

\subsubsection{Residual Standard Deviation $\sigma$: Gradual Convergence}

The residual standard deviation $\sigma$ follows an interesting trajectory: its
posterior mean rises from 652.3 in Block 1 to 1{,}497.1 in Block 4, converging
closely to the full-data estimate of approximately 1{,}500. Block 1's low estimate
reflects the reduced variability within a single semi-annual window compared to
the full two-year span. The moment-matched prior for Block 2 is
$\operatorname{Half-Normal}(652)$, which places most prior mass below the true
value but does not prevent the likelihood from pushing $\sigma$ upward toward
1{,}127. By Block 4, the data have effectively overridden the prior and
$\sigma$ has converged to the correct magnitude.

\subsubsection{Implications for Bayesian Learning}

The sequential analysis demonstrates two complementary aspects of Bayesian
learning. For a dominant signal such as the temperature effect, sequential
updating reliably narrows uncertainty and stabilises estimates even with
$n = 90$ per block. For a weaker signal such as the humidity effect, the same
procedure reveals the limits of small-sample inference: local estimates are
dominated by noise and seasonal confounding, and the propagation of moment-matched
priors can entrench early estimation errors rather than correcting them. This
asymmetry highlights an important practical consideration: sequential Bayesian
updating is most reliable when the parameter of interest is well-identified
within each individual segment.
